\chapter{Your first binding}
\label{ch:first}

This chapter goes from an unmodified header to an importable Python module and
a loadable Node addon. It is deliberately complete: every command, in order,
with what each one produces.

\section{Prerequisites}

\begin{itemize}
  \item A \textbf{clang-p2996} build. Rosetta looks for it at
    \code{\$ENV\{HOME\}/devs/c++/clang-p2996/build} by default; anywhere else
    is fine, you just say where (Section~\ref{sec:toolchain-paths}).
  \item \textbf{CMake} 3.28 or later, and Ninja or Make.
  \item For the Python target, \code{pip install pybind11}. For Node,
    \code{npm}. For WebAssembly, an activated emsdk.
\end{itemize}

The reflection flags rosetta's generated projects use, should you need to know
them: \code{-freflection -freflection-latest -fexperimental-library}, plus
\code{-fannotation-attributes} for anything using annotations.

\begin{note}
Only the \emph{generation host} needs the fork. Everything under
\code{bindings/} builds with a stock compiler --- with the two documented
exceptions, \lang{rest} and \lang{json}.
\end{note}

\section{Build the tool, once}

\begin{lstlisting}[style=sh]
cmake -G Ninja -S tools/rosetta_gen -B tools/rosetta_gen/build
cmake --build tools/rosetta_gen/build
\end{lstlisting}

The binary always lands in \code{<repo>/bin/rosetta\_gen} --- not in the build
tree. This step needs no reflection: \code{rosetta\_gen} is a JSON reader and a
text templater, and builds with an off-the-shelf C++17 compiler.

\section{The library --- leave it alone}

\begin{lstlisting}[style=cpp]
// my_lib/person.h
#pragma once
#include <string>

struct Person {
    std::string name;
    int         age = 0;
    std::string id;

    Person() = default;
    Person(std::string n, int a) : name(std::move(n)), age(a) {}

    std::string greet(const std::string &salutation) const;
};
\end{lstlisting}

Note what is \emph{not} there: no macro, no registration call, no rosetta
include, no export annotation. That is the whole point.

\section{The manifest}

Put it wherever you like --- every relative path inside resolves from the
manifest's own directory. Here, \code{my\_lib/rosetta/manifest.json}:

\begin{lstlisting}[style=json]
{
  "user_include": "..",
  "rosetta_include": "/path/to/rosetta/include",
  "generator_name": "person_gen",
  "module_name": "person",

  "targets": [
    { "lang": "python", "name": "person_py" },
    { "lang": "node",   "name": "person_js" },
    "typescript",
    "markdown"
  ],

  "classes": [
    { "name": "Person", "header": "person.h" }
  ]
}
\end{lstlisting}

Four required ideas and nothing else:

\begin{description}[leftmargin=0cm,style=nextline,font=\normalfont\code]
  \item[user\_include] where your headers live. One directory, or an array of
    them.
  \item[rosetta\_include] where rosetta's \code{include/} lives.
  \item[targets] which backends to emit. A bare string uses
    \field{module\_name}; the object form sets a per-target module name.
  \item[classes] what to bind. \field{header} is required; \field{name}
    defaults to the header's stem.
\end{description}

\field{generator\_name} is optional --- it names the emitted driver and
defaults to the manifest's folder name.

\begin{note}
Do not want to write this by hand? \code{rosetta\_gen --init} writes a
fully-commented starter manifest, and \code{rosetta\_gen --init manifest.json
./src} pre-fills it from a scan of your source tree.
See Section~\ref{sec:init}.
\end{note}

\section{Generate and build --- the one-command version}

\begin{lstlisting}[style=sh]
/path/to/rosetta/bin/rosetta_gen --build manifest.json
\end{lstlisting}

That emits the generator project, builds it (this is where reflection runs),
runs it, and then compiles every backend. The run ends with a summary:

\begin{lstlisting}[style=plain]
== summary
   python      OK
   node        OK
   typescript          OK (nothing to compile)
   markdown            OK (nothing to compile)
\end{lstlisting}

A backend whose toolchain is missing is skipped with a note rather than failing
the run:

\begin{lstlisting}[style=plain]
   wasm        SKIPPED (emcc not found -- activate emsdk)
\end{lstlisting}

\section{The same thing, one step at a time}

Worth doing once, because it makes the pipeline concrete.

\begin{lstlisting}[style=sh]
# 1. manifest -> generator project (no reflection yet)
rosetta_gen manifest.json gen

# 2. build the generator -- REFLECTION HAPPENS HERE
cmake -S gen -B gen/build && cmake --build gen/build

# 3. run it -- one self-contained project per target
./generator bindings
\end{lstlisting}

After step~1, look inside \code{gen/}. There are three files:

\begin{center}
\small
\begin{tabular}{@{}lL{9.2cm}@{}}
\toprule
\code{person\_gen.cpp} & a \code{main()} that calls
  \code{rosetta::generate<Person>(\ldots)} --- the driver \\
\code{bindings.h} & your headers included, plus any baked-in out-of-line
  annotations \\
\code{CMakeLists.txt} & knows where clang-p2996 is and which flags it needs \\
\bottomrule
\end{tabular}
\end{center}

After step~3:

\begin{lstlisting}[style=plain]
bindings/
  coverage.json
  python/      auto_pybind.cpp  CMakeLists.txt  pyproject.toml  make_wheel.py
  node/        auto_napi.cpp    CMakeLists.txt  package.json
  typescript/  person.d.ts
  markdown/    person.md
\end{lstlisting}

Open \code{auto\_pybind.cpp}. It is ordinary pybind11 --- explicit
\code{.def()} calls, one per member, with no reflection anywhere. That file is
shippable: it compiles on any machine with pybind11 and a normal C++
compiler.

\section{Build and use it}

\subsection*{Python}

\begin{lstlisting}[style=sh]
cd bindings/python
cmake -G Ninja -B build && cmake --build build
\end{lstlisting}

\begin{lstlisting}[style=py]
import person_py

p = person_py.Person("Alice", 30)
p.name = "Alice B."
print(p.greet("Hello"))
print(p.age)
\end{lstlisting}

\subsection*{Node}

\begin{lstlisting}[style=sh]
cd bindings/node
npm install && npx cmake-js compile
\end{lstlisting}

\begin{lstlisting}[style=js]
const m = require('./build/Release/person_js.node');
const p = new m.Person('Alice', 30);
console.log(p.greet('Hello'));
\end{lstlisting}

\subsection*{The two you did not have to build}

\code{bindings/typescript/person.d.ts} declares the Node addon's surface, so an
editor autocompletes it. \code{bindings/markdown/person.md} is an API reference
page. Neither compiles anything --- they are text, produced directly.

\section{Read \code{coverage.json} once}

Before moving on, open \code{bindings/coverage.json}. Even for a three-field
class it tells you something worth having: exactly which members reached each
backend and which did not, with a reason for every drop.

\begin{lstlisting}[style=json]
{
  "rosetta_coverage": 1,
  "reflection": [],
  "targets": [
    { "target": "python", "bound": 5, "skipped": 0, "classes": [ ... ] }
  ]
}
\end{lstlisting}

On a real library it will not be empty, and the habit of checking it is the
difference between knowing your binding surface and guessing at it.
Chapter~\ref{ch:coverage} is about reading it properly.

\section{Enriching, if you want to}

Everything above used reflection alone: names match the C++ identifiers, types
map automatically, and you get no docstrings, no validation and no read-only
semantics. Adding them is opt-in, member by member:

\begin{lstlisting}[style=cpp]
#include <rosetta/annotations.h>

struct Person {
    [[= rosetta::doc{"display name"}]]
    std::string name;

    [[= rosetta::doc{"age in years"}, = rosetta::range{0.0, 150.0}]]
    int age = 0;

    [[= rosetta::doc{"server-assigned id"}, = rosetta::readonly{}]]
    std::string id;
};
\end{lstlisting}

Now \code{help(person\_py.Person.name)} says something, assigning \code{200} to
\code{age} raises, and \code{id} has no setter.

\begin{warn}
Writing annotations inline makes the header C++26-only --- every translation
unit that includes it, including the generated binding, now needs the fork.
If you want to keep the header plain C++ (and you usually do), use the JSON
side-car instead: same metadata, same result, one manifest field.
See Section~\ref{sec:outofline}.
\end{warn}

\section{Cleaning up}

\begin{lstlisting}[style=sh]
rosetta_gen --clean manifest.json
\end{lstlisting}

removes \code{gen/}, the generator binary, and the per-backend folders under
\code{bindings/} \emph{that this manifest declares}. Anything else you parked
there survives, and a folder whose \code{CMakeLists.txt} lacks the
\code{Generated by rosetta} stamp is never touched. The manifest itself is
never removed.

\section{What to read next}

You now have a working binding. Part~II is the manifest in full --- and the
next chapter starts with the rules that apply to all of it.
